% frontmatter-2B-v3-20261124.tex --- main paper B (2A + 2B merged) frontmatter.
% NEW FILE, 2026-11-24. Supersedes frontmatter-2B-v2.tex for the assembled draft;
% v2 is left untouched (only-add discipline).
% Changes vs v2: title drops "$e$"; abstract merges the 2A and 2B claims and
% demotes the convergent denominators of $e$ to a flagship corollary; the
% unified structural-alphabet placeholder is retained as an explicit open item.
\title{Structural alphabets, exact Hausdorff gauges, and dimension phase transitions}

\begin{document}
\maketitle

\begin{abstract}
For polynomial-growth recurrences $Q_{j+1}=(cj^r+d)Q_j+Q_{j-1}$ and for
arbitrary growing-quotient recurrences $Q_{j+1}=A_jQ_j+Q_{j-1}$ with
$A_j\to\infty$, we determine the full spectrum of the critical exponent
$\Lam(x)=\limsup_j j^r\|Q_jx\|$, the zero-level lattice, the exact critical
gauges, the positive-level three-way splitting, and the dimension phase
transition. For $A_j=4j+2$, the convergent denominators of $\ee$, the spectrum
is $\frac14\Z_{\ge0}\cup\{\infty\}$ with zero level $L=\Z+\Z\frac{\ee+1}{2}$,
every finite level has Hausdorff dimension $0$, and the gauge
$h_k(1/Q_n)=(2k+1)^{-n}$ is exact. For finite-state defect constraints the
critical gauge is governed by the largest Perron root among marked components
supporting infinitely many extreme defects. All results are machine-checked in
Lean~4 using only the base axioms \texttt{propext}, \texttt{Classical.choice}
and \texttt{Quot.sound}, except the finite-dimensional Perron--Frobenius
eigenvector existence input, passed explicitly as a hypothesis. This work does
not resolve the open problem of the irrationality of $e+\pi$.
\end{abstract}

\section*{Honesty statement and formal verification}

The Lean~4 audit subset of this paper --- the modules covering every theorem
stated above, with names transcribed verbatim from the corresponding sections
of the full manuscript --- is: \texttt{DeltaSign.lean},
\texttt{EpiFormal/Classification.lean}, \texttt{EpiFormal/SpectralGap.lean},
\texttt{EpiFormal/FullSpectrumDiscrete.lean}, \texttt{EpiFormal/DimHCore.lean},
\texttt{EpiFormal/DimHFinite.lean}, \texttt{EpiFormal/GaugeExact.lean},
\texttt{QuarterPointFinal.lean}, \texttt{EpiFormal/UnweightedCorrection.lean},
\texttt{WeightedFiber.lean}, \texttt{WeightedDimHCore.lean},
\texttt{WeightedUpper.lean}, \texttt{WeightedFinal.lean},
\texttt{EpiFormal/WeightedExactDim.lean}, \texttt{EpiFormal/GaugeNull.lean},
\texttt{EpiFormal/GaugeCircle.lean}, \texttt{GaugeLiminfComplete.lean},
\texttt{EpiFormal/TrichotomyAndPolyExtension.lean},
\texttt{EpiFormal/FiniteStateComplete.lean} (building on
\texttt{FiniteStateGauge.lean}, \texttt{FiniteStateTail.lean}, and
\texttt{FiniteStatePositive.lean}),
\texttt{EpiFormal/PolyGaugeAndPhaseTransition.lean},
\texttt{EpiFormal/PolyWeightedUpper.lean},
\texttt{EpiFormal/PolyWeightedLower.lean},
\texttt{EpiFormal/GeneralGrowingRec.lean}, and
\texttt{EpiFormal/GeneralGrowingDim.lean}. The full manuscript reports these
modules as kernel-checked at Layer~A with only the three base axioms and zero
\texttt{sorryAx}; they are audited jointly with the companion papers' modules
in \texttt{Audit.lean}, whose declaration counts are project-wide, not
per-paper. The finite-dimensional linear-algebraic Perron--Frobenius
eigenvector existence input on marked strongly connected components, used by
the finite-state tail theorem, is passed explicitly as a hypothesis and is not
derived inside the Lean kernel. The unified structural-alphabet theorem is not
yet proved; it appears in this paper only as an explicitly marked open item.

\emph{The full version containing all proofs is available as a companion manuscript.}
